\documentclass[10pt]{article}

% Specify the margins and text size.
\setlength{\textwidth}{6.4in}
\setlength{\textheight}{9.6in}
\setlength{\oddsidemargin}{0pt}
\setlength{\evensidemargin}{0pt}
\setlength{\topmargin}{0pt}
\setlength{\hoffset}{.05in}
\setlength{\voffset}{-1in}

\setlength{\parskip}{5pt}
\setlength{\parindent}{0pt}

% Load some fonts and symbol packages
\usepackage{latexsym}
\usepackage{pifont}       % contains 'star' symbol for counterinsurgency handbook title
\usepackage{yfonts} 
\usepackage{amsmath}
\usepackage{amsfonts}

\usepackage{graphicx}     % actually, this is loaded by pstricks
\usepackage[T1]{fontenc}
\usepackage{ifthen}
\usepackage{pstricks,pst-func,pst-text,pst-node,multido,pst-plot,calc,pst-3dplot}
%\usepackage[all]{xy}
%\usepackage{animate}

% The hyperref package inserts links into the pdf.
\definecolor{MyLinkColor}{rgb}{.1,.2,1}
\definecolor{MyCiteColor}{rgb}{.1,1,.2}
\definecolor{MyURLColor}{rgb}{.4,.4,.4}
\usepackage[backref=true,pagebackref=false,hyperindex,colorlinks=true,
  linkcolor=MyLinkColor,urlcolor=MyURLColor]{hyperref}


% The tweaklist package is something I found on the web.  It provides a simple interface
% for making changes to spacing used in the itemize and enumerate environments.  Comment
% this out if you don't care to use tweaklist.
\usepackage{tweaklist}
\renewcommand{\itemhook}{\setlength{\parskip}{2pt}\setlength{\parsep}%
{1pt}\setlength{\topsep}{0pt}\setlength{\itemsep}{0pt}}

\newcommand{\N}{\ensuremath\mathbb{N}}
\newcommand{\Q}{\ensuremath\mathbb{Q}}
\newcommand{\R}{\ensuremath\mathbb{R}}
\newcommand{\C}{\ensuremath\mathbb{C}}
\newcommand{\z}{\hphantom{0}}

\begin{document}
\pagestyle{empty}

\href{http://www.su.edu}{\includegraphics[height=2.2cm]{sulogo.eps}}
\vspace{-1.69cm}

{\small
\begin{tabular}{cl}
\hspace{5in} & Math 207\\
& Intro to Statistics
%                & Worksheet\\
%                & %1 February 2012
\end{tabular}
}
\medskip

{\setlength{\baselineskip}{1.05\baselineskip}
\begin{center}
\textbf{\large  Worksheet:  Hypothesis Testing}
\end{center}

1. \textbf{Areas Under the Normal Curve}.
Sketch the following regions under the normal curve and estimate their areas.\vspace{-3pt}
\begin{center}
\begin{tabular}{lcl}
a) to the left of $z=0$ &\hspace{2in} & b) between $z=-1$ and $z=2$\\[.7in]
c) between $z=-2$ and $z=1$ &\hspace{2in} & d) to the left of $z=3$ \\[.7in]
e) to the right of $z=2$ &\hspace{2in} & f) $z>5$\\[.7in]
\end{tabular}
\end{center}

2. \textbf{Standard Units}. A list of numbers has $\mbox{mean}=4$ and $\mbox{SD}=2$.

\hspace{10pt} a) Convert $x=6$ to standard units.
\vspace{.5in}

\hspace{10pt} b) What value $x$ as value $-2$ in standard units.
\vspace{.5in}


3. \textbf{Central Limit Theorem}.
Would taking the average of 25 measurements likely divide the error by a factor of 5, 10 or 25?
\vspace{.3in}

4. \textbf{$p$-Values}.

\hspace{10pt} a) Which of the following $p$-values is best for the
null hypothesis? Explain.
\[0.1\%\hspace{.3in} 3\%\hspace{.3in} 17\%\hspace{.3in} 32\%\]
\vspace{.1in}

\hspace{10pt} b) Repeat a) for the alternative hypothesis.
\vspace{.3in}


\hspace{10pt} c) True or false:  The $p$-value depends on the data.
\vfill


\hspace{10pt} d) True or false: If the $p$-value is 5\%, there is only a 5\% chance that
the null hypothesis is correct.
\eject

5. \textbf{Hypothesis Tests}.
For each question,  (1) state the test you would use, (2) state the null and alternative hypotheses,
(3) compute the test statistic and $p$-value (approximately) and (4) make a conclusion.

\hspace{10pt} a)  You flip a coin 100 times and it lands heads 51\% of the time.  Is that too many heads?
\vspace{1.5in}

\hspace{10pt} b) Repeat a) if  $n=10,000$.  
\vspace{1.5in}

\hspace{10pt} c) A coin lands heads four times in a row.   Is that too many heads?
\vspace{1.5in}

\hspace{10pt} d)   A local pizza shop advertises that
\lq\lq 75\% of Shenandoah students enjoy our pizza every week!\rq\rq\
You take a  random sample of 100 SU students and find that 50 of them ordered 
 pizza from that shop last week.  Do the data support the advertisement?
\vspace{1.5in}

\hspace{10pt} e)
In 1970, 59\% of college freshmen thought that capital punishment
should be abolished;  by 2005, the percentage had dropped to 35\%.
Is the difference real, or can it be explained by chance? You may assume that
the percentages are based on two independent simple random samples, each of
size 1,000.
\vfill
\eject

\begin{center}
\textbf{\large Solutions}
\end{center}

{\small

1. a) $50\%$ \hspace{3pt}
   b) $(95\%/2) + (68\%/2) = 81.5\%$ \hspace{3pt}
   c) same as b) \hspace{3pt}
   d) $(99.7\% + (0.3\%/2) = 99.85\%$ \hspace{3pt}
   e) about 0\%

2. a) $z=(6-4)/2 = 2/2 = 1$\hspace{3pt} b) Solve $-2 = (x-4)/2$ to get $x=0$.

3. $\mbox{SE}_{\mbox{\scriptsize av}} = \mbox{SD}/\sqrt{n} = \mbox{SD}/\sqrt{25} = \mbox{SD}/5$.  So, the answer is 5.

4. a) 32\%.  If the null hypothesis is true, there is a 32\% chance of seeing data as extreme or more extreme
than was observed.  b) $0.1\%$.  If the null were true then there would be very little chance of seeing
such extreme data.  c) True.  d) False.  We assume that the null is true in order to compute the $p$-value.

5. a) $z$-test.   Null hypothesis:  the difference between expected (50\%) and observed (51\%) is due to chance.
Alternative hypothesis:  No, it's not.  The box model has a 0 and 1 in it.  It's mean is $0.5$
and $\mbox{SD}$ is $0.5$.
The SE is $0.5/\sqrt{100} = 5\%$.  So, the $z$-score is $(51-50)/5 \approx 0.2$  So, the $p$-value is
close to 50\% (it's 42.1\%).  This supports the null.

b) $z$-test.  Null, alt and box are the same.  $\mbox{SE}$ is 0.5\% and $z=(51-50)/0.5 =2$.  That gives a $p$-value
of about 2.5\% which supports the alternative.

c) $t$-test since $n$ is small.  Null, alt and box are the same as above.
$\mbox{SE}$ is $0.5/\sqrt{4}=0.25 = 25\%$.  So, $z=(0-50)/25 \approx -2$.
If we use a $t$-distribution with 3 degrees of freedom, we get a $p$-value of about 7\% which supports the null.
Note that if we had used the $z$-test, the $p$-value would be 2.5\% (which would support the alternative).

d) $z$-test.  Null hypothesis:  The difference between expected (75\%) and observed (50\%) is due to chance.
Alternative hypothesis:  No, it's not.  Box:  3 ones and a zero.  $\mbox{AV}_{\mbox{\scriptsize box}} = 0.75$,
$\mbox{SD}_{\mbox{\scriptsize box}} = 0.43$,
$\mbox{SE} = 0.43/\sqrt{100} = 0.043 = 4.3\%$.  The test statistic is $z=(50-75)/4.3 < -5$ so the $p$-value is
about 0\%.  This supports the alternative.

e) Two-sample $z$-test.  Null hypothesis:  The observed difference (50\% and 55\%) is due to chance.
Alternative hypothesis:  No, it's not.  For the 1970s box, the SD is $\sqrt{0.59 * 0.41}\approx 0.49$
so the SE is $0.49/\sqrt{1000} \approx 1.5\%$.  For the 2005 box, the SD is
$\sqrt{0.35* 0. 65} \approx 0.48$ so the SE  is $0.48/\sqrt{1000} \approx 1.5\%$.
The SE for the difference is $\sqrt{1.5^2 + 1.5^2}\approx 2.1\%$.  So, the test statistic
is $z = (59 - 35)/2.1 = 11.4$.  This gives a $p$-value close to 0 which supports the alternative hypothesis.}
\vfill

\begin{center}
\textbf{\large Solutions}
\end{center}

{\small

1. a) $50\%$ \hspace{3pt}
   b) $(95\%/2) + (68\%/2) = 81.5\%$ \hspace{3pt}
   c) same as b) \hspace{3pt}
   d) $(99.7\% + (0.3\%/2) = 99.85\%$ \hspace{3pt}
   e) about 0\%

2. a) $z=(6-4)/2 = 2/2 = 1$\hspace{3pt} b) Solve $-2 = (x-4)/2$ to get $x=0$.

3. $\mbox{SE}_{\mbox{\scriptsize av}} = \mbox{SD}/\sqrt{n} = \mbox{SD}/\sqrt{25} = \mbox{SD}/5$.  So, the answer is 5.

4. a) 32\%.  If the null hypothesis is true, there is a 32\% chance of seeing data as extreme or more extreme
than was observed.  b) $0.1\%$.  If the null were true then there would be very little chance of seeing
such extreme data.  c) True.  d) False.  We assume that the null is true in order to compute the $p$-value.

5. a) $z$-test.   Null hypothesis:  the difference between expected (50\%) and observed (51\%) is due to chance.
Alternative hypothesis:  No, it's not.  The box model has a 0 and 1 in it.  It's mean is $0.5$
and $\mbox{SD}$ is $0.5$.
The SE is $0.5/\sqrt{100} = 5\%$.  So, the $z$-score is $(51-50)/5 \approx 0.2$  So, the $p$-value is
close to 50\% (it's 42.1\%).  This supports the null.

b) $z$-test.  Null, alt and box are the same.  $\mbox{SE}$ is 0.5\% and $z=(51-50)/0.5 =2$.  That gives a $p$-value
of about 2.5\% which supports the alternative.

c) $t$-test since $n$ is small.  Null, alt and box are the same as above.
$\mbox{SE}$ is $0.5/\sqrt{4}=0.25 = 25\%$.  So, $z=(0-50)/25 \approx -2$.
If we use a $t$-distribution with 3 degrees of freedom, we get a $p$-value of about 7\% which supports the null.
Note that if we had used the $z$-test, the $p$-value would be 2.5\% (which would support the alternative).

d) $z$-test.  Null hypothesis:  The difference between expected (75\%) and observed (50\%) is due to chance.
Alternative hypothesis:  No, it's not.  Box:  3 ones and a zero.  $\mbox{AV}_{\mbox{\scriptsize box}} = 0.75$,
$\mbox{SD}_{\mbox{\scriptsize box}} = 0.43$,
$\mbox{SE} = 0.43/\sqrt{100} = 0.043 = 4.3\%$.  The test statistic is $z=(50-75)/4.3 < -5$ so the $p$-value is
about 0\%.  This supports the alternative.

e) Two-sample $z$-test.  Null hypothesis:  The observed difference (50\% and 55\%) is due to chance.
Alternative hypothesis:  No, it's not.  For the 1970s box, the SD is $\sqrt{0.59 * 0.41}\approx 0.49$
so the SE is $0.49/\sqrt{1000} \approx 1.5\%$.  For the 2005 box, the SD is
$\sqrt{0.35* 0. 65} \approx 0.48$ so the SE  is $0.48/\sqrt{1000} \approx 1.5\%$.
The SE for the difference is $\sqrt{1.5^2 + 1.5^2}\approx 2.1\%$.  So, the test statistic
is $z = (59 - 35)/2.1 = 11.4$.  This gives a $p$-value close to 0 which supports the alternative hypothesis.}


\eject

\end{document}
